% Hand-written discussion for make_report.py. It interprets the numbers of one evaluation run; make_report.py
% flags it when the run it is built with is a different one.
% written for run: 2026-10-14T21:27:23
\section{Discussion}

\subsection{COLMAP gives the most accurate trajectory overall}
COLMAP posed all 344 frames (all 688 images in one model) with a position error of 2.43~mm RMS and 5.4~mm at
worst over 848~mm of travel, an orientation error of 0.49$^\circ$, and a best-fit scale of 0.999. It solves every
frame together and holds the calibrated 150~mm baseline fixed, so its error is spread evenly along the pass rather
than accumulating (Figure~\ref{fig:fig_components}), and its scale is metric without any tag. It needs no depth,
which makes it the cheapest method here once the depth needed by ICP is counted (12 minutes against 19--28).

\subsection{ICP is better locally but drifts}
ICP on FoundationStereo depth with a local map comes close: 2.89~mm RMS, 7.1~mm worst. Locally it is the better of
the two. Its motion over 50~mm is wrong by 2.5\% (median) against COLMAP's 4.1\%, and its orientation error is
lower (0.38$^\circ$). Over the whole pass the small errors add up: its scale is 1\% short (0.990), which shows as
the steady slope of the along-track error in Figure~\ref{fig:fig_components}. The whole difference between the two
methods is in that slope: one is a chain, the other a joint solution. A pass several times longer would separate
them further, unless ICP is anchored by something like COLMAP or loop closure.

\subsection{The depth decides ICP}
On SGBM depth every ICP variant is about twice as wrong as on FoundationStereo depth (5.2 against 2.9~mm with the
local map; 10.1 against 4.4~mm frame-to-frame). SGBM keeps about half as many points per frame and leaves the
dark, weakly textured prints and much of the tray empty, which are exactly the surfaces that face along the track
and hold ICP against sliding. Its median translation constraint is 0.05 against 0.13 for FoundationStereo.

\subsection{A local map halves ICP's error; colour adds nothing}
Registering onto the last ten frames instead of only the previous one halves the error with either depth source.
A larger overlap gives more surface facing along the track, and one noisy frame pulls the chain less. Coloured ICP
changed nothing: within 0.02~mm with FoundationStereo and 0.1~mm with SGBM, at three to four times the run time.
At Open3D's default weighting (geometry 0.968, colour 0.032), and on a grey image whose texture is mostly dark
prints on a white board, the colour term barely moves the solution.

\subsection{What holds the methods in place, and what that means for ice}
Figures~\ref{fig:det_a} and~\ref{fig:det_b} show where the grip comes from. COLMAP's features sit on edges and
speckle: the brick texture, the tray rim and handle, the outlines of the objects, the seams between the board
panels and specks of debris on the white board. The flat white board and the inside of the dark prints give
almost nothing. Where the bricks are in view they are the single largest source (only 38\% of the triangulated
features of frame 0610 lie on the board). ICP's grip against the slide comes from the pyramid faces, the object
walls and the curved tray rim. The board itself, flat and parallel to the motion, only holds ICP in depth.

Both methods therefore lean on structure that a bare ice surface may not have. These results are a best case
for both. The translation constraint is the number to watch on ice: here it is 0.05--0.13, already low. A
smooth, flat sheet of ice would push it towards zero, and point-to-plane ICP would then drift along the track
with nothing to stop it. Ice with ridges, cracks, frost texture or embedded particles could serve both methods,
and only footage of the real surface can say which.

\subsection{Limitations}
\begin{itemize}
\item One recording, one pass, one distance (s30), every 5th frame. The ICP fixes of Section~\ref{sec:dev} were
made on this recording.
\item The tag poses carry their own noise (0.28~mm median layout residual), and 17 of the 344 frames had no tag
found, so a tag may stay visible in them.
\item Only the poses are scored. Whether a 2--3~mm pose error matters depends on what it does to the fused
surface, which this experiment does not yet measure.
\end{itemize}

\subsection{Next steps}
\begin{enumerate}
\item Fuse the per-frame depth with each pose set and score the pyramid against CAD, to see whether the 2.4 against
2.9~mm difference in pose shows in the surface.
\item Repeat on a second recording (another distance, e.g.\ s40) with the settings frozen, to give the ICP rows an
untuned test.
\item Hide everything except the board (a mask of the surroundings) to see how each method fares with only what
ice-like surfaces offer.
\item Run every estimator on a clip over real ice: no reference, but registration rate and translation
constraint still show whether they hold.
\item Try COLMAP poses refined by ICP, which combines COLMAP's global consistency with ICP's local accuracy.
\end{enumerate}
